% Assignment 1 -- GAN Playground
% Deep Neural Networks for Computer Graphics
%
% This file is self-contained and compiles cleanly on Overleaf with
% the default pdfLaTeX engine.

\documentclass[11pt,a4paper]{article}

% --- packages -------------------------------------------------------------
\usepackage[utf8]{inputenc}
\usepackage[T1]{fontenc}
\usepackage[margin=1in]{geometry}
\usepackage{amsmath,amssymb,amsthm}
\usepackage{graphicx}
\usepackage{booktabs}
\usepackage{hyperref}
\usepackage{xcolor}
\usepackage{listings}
\usepackage{enumitem}
\usepackage{caption}

\hypersetup{
    colorlinks=true,
    linkcolor=blue!60!black,
    urlcolor=blue!60!black,
    citecolor=blue!60!black,
}

% --- listings: python ----------------------------------------------------
\definecolor{codebg}{rgb}{0.97,0.97,0.97}
\definecolor{codekw}{rgb}{0.13,0.29,0.53}
\definecolor{codestr}{rgb}{0.50,0.30,0.10}
\definecolor{codecmt}{rgb}{0.40,0.55,0.40}

\lstdefinestyle{py}{
    language=Python,
    backgroundcolor=\color{codebg},
    basicstyle=\ttfamily\small,
    keywordstyle=\color{codekw}\bfseries,
    stringstyle=\color{codestr},
    commentstyle=\color{codecmt}\itshape,
    showstringspaces=false,
    breaklines=true,
    frame=single,
    framesep=4pt,
    rulecolor=\color{black!15},
    numbers=left,
    numberstyle=\tiny\color{black!50},
    numbersep=8pt,
    columns=fullflexible,
}
\lstset{style=py}

% --- title ---------------------------------------------------------------
\title{Assignment~1 -- GAN Playground\\
       \large Deep Neural Networks for Computer Graphics}
\author{Roi Moshe}
\date{\today}

% =========================================================================
\begin{document}
\maketitle

% =========================================================================
\section*{Part 1 -- Deep Convolutional GAN}

% -------------------------------------------------------------------------
\subsection*{Task 1.1 -- Data Augmentation}

The \texttt{vanilla} branch of \texttt{get\_data\_loader} composes the
following pipeline (in order): resize to a slightly larger canvas, take a
random square crop back to \texttt{opts.image\_size}, randomly flip
horizontally, convert to a tensor in $[0,1]$, then normalise to $[-1,1]$
so the output range matches the generator's \texttt{tanh}.

\begin{lstlisting}
load_size = int(1.1 * opts.image_size)
osize = [load_size, load_size]
train_transform = transforms.Compose([
    transforms.Resize(osize, Image.BICUBIC),
    transforms.RandomCrop(opts.image_size),
    transforms.RandomHorizontalFlip(),
    transforms.ToTensor(),
    transforms.Normalize((0.5, 0.5, 0.5), (0.5, 0.5, 0.5)),
])
\end{lstlisting}

% -------------------------------------------------------------------------
\subsection*{Task 1.2a -- Padding Derivation}

For a 2-D convolution with input spatial size $H_{\text{in}}$, kernel
size $K$, stride $S$, and padding $P$, the output size is
%
\begin{equation}
    H_{\text{out}}
    \;=\;
    \left\lfloor
    \frac{H_{\text{in}} + 2P - K}{S}
    \right\rfloor + 1.
    \label{eq:conv-out}
\end{equation}
%
The discriminator's first four convolutions must downsample by exactly a
factor of two, i.e. $H_{\text{out}} = H_{\text{in}}/2$, with the prescribed
$K = 4$ and $S = 2$. Substituting into~\eqref{eq:conv-out} (and noting
that $H_{\text{in}}$ is even so the floor is exact):
%
\begin{align*}
    \frac{H_{\text{in}}}{2}
        &= \frac{H_{\text{in}} + 2P - 4}{2} + 1 \\[2pt]
    H_{\text{in}}
        &= H_{\text{in}} + 2P - 4 + 2 \\[2pt]
    0   &= 2P - 2 \\[2pt]
    \boxed{P = 1.}
\end{align*}
%
Sanity check on layer~1: $H_{\text{in}}=64$ gives
$\lfloor (64 + 2 - 4)/2 \rfloor + 1 = 31 + 1 = 32$, halving the spatial
extent as required. The same holds for the remaining three downsampling
layers.

% -------------------------------------------------------------------------
\subsection*{Task 1.2b -- DCDiscriminator}

Using $P=1$ for layers~1--4 and the helper \texttt{conv} from
\texttt{models.py}, the discriminator follows the channel/spatial schedule
in Table~\ref{tab:disc}. The final \texttt{conv5} collapses the
$4{\times}4$ feature map to a single scalar with $K=4,\,S=1,\,P=0$ (and
no normalisation or activation, since the loss is least-squares on the
raw logit).

\begin{table}[h]
    \centering
    \begin{tabular}{lcccccc}
        \toprule
        Layer & $C_{\text{in}}\!\to\!C_{\text{out}}$ & $K$ & $S$ & $P$ & Norm & Activation \\
        \midrule
        conv1 & $3 \to 32$    & 4 & 2 & 1 & InstanceNorm & ReLU \\
        conv2 & $32 \to 64$   & 4 & 2 & 1 & InstanceNorm & ReLU \\
        conv3 & $64 \to 128$  & 4 & 2 & 1 & InstanceNorm & ReLU \\
        conv4 & $128 \to 256$ & 4 & 2 & 1 & InstanceNorm & ReLU \\
        conv5 & $256 \to 1$   & 4 & 1 & 0 & --           & --    \\
        \bottomrule
    \end{tabular}
    \caption{DCDiscriminator layer schedule. Output spatial sizes follow
             the doubling-down sequence $64 \to 32 \to 16 \to 8 \to 4 \to 1$.}
    \label{tab:disc}
\end{table}

% =========================================================================
% Future tasks will be appended below as the assignment progresses.
% =========================================================================

\end{document}
